\documentclass[11pt]{article}
\usepackage[T1]{fontenc}
\usepackage[utf8]{inputenc}
\usepackage{lmodern}
\usepackage[english]{babel}
\usepackage{amsmath,amssymb,mathtools,array,booktabs,pdflscape,adjustbox,diagbox}
\usepackage{geometry}
\usepackage{microtype}
\geometry{a4paper,margin=1.45cm}
\setlength{\parindent}{1.25em}
\setlength{\parskip}{0pt}
\makeatletter
\renewcommand\footnotesize{\@setfontsize\footnotesize{7.5}{8.5}\selectfont}
\makeatother
\providecommand{\widebar}[1]{\overline{#1}}
\providecommand{\A}{\Delta}
\providecommand{\D}{\Delta}
\allowdisplaybreaks
\setlength{\emergencystretch}{10em}
\sloppy
\hbadness=10000
\vbadness=10000
\hfuzz=2pt
\vfuzz=10pt
\raggedbottom
\providecommand{\R}{\varrho}
\providecommand{\hata}{\widehat{a}}
\providecommand{\hatv}{\widehat{v}}
\providecommand{\modu}[1]{\;\operatorname{mod} #1}
\providecommand{\mcell}[1]{\ensuremath{\begin{gathered}#1\end{gathered}}}
\providecommand{\pair}[2]{(#1\,|\,#2)}
\providecommand{\eps}{\varepsilon}
\providecommand{\rhoidx}{\varrho}

% Batch 20 macros
\providecommand{\dd}{\mathrm d}
\providecommand{\Div}{\operatorname{Div}}
\providecommand{\G}{\mathfrak G}
\providecommand{\bd}{\bar{\delta}}
\providecommand{\p}{\partial}


% Batch 21 ideal/function-field macros
\providecommand{\frS}{\mathfrak S}
\providecommand{\frp}{\mathfrak p}
\providecommand{\fra}{\mathfrak a}
\providecommand{\frb}{\mathfrak b}
\providecommand{\frc}{\mathfrak c}
\providecommand{\frf}{\mathfrak f}
\providecommand{\frr}{\mathfrak r}
\providecommand{\frm}{\mathfrak m}
\providecommand{\frD}{\mathfrak D}
\providecommand{\calA}{\mathcal A}
\providecommand{\calB}{\mathcal B}
\providecommand{\calN}{\mathcal N}


\providecommand{\p}{\partial}
\providecommand{\dd}{\mathrm d}
\providecommand{\modu}[1]{\;\operatorname{mod} #1}
\providecommand{\mM}{\mathfrak{M}}
\providecommand{\mN}{\mathfrak{N}}
\providecommand{\mL}{\mathfrak{L}}
\providecommand{\mG}{\mathfrak{G}}
\providecommand{\mH}{\mathfrak{H}}
\providecommand{\mH}{\mathfrak{H}}
\providecommand{\mA}{\mathfrak{A}}
\providecommand{\mB}{\mathfrak{B}}
\providecommand{\mX}{\mathfrak{X}}
\providecommand{\mY}{\mathfrak{Y}}
\providecommand{\ma}{\mathfrak{a}}
\providecommand{\mb}{\mathfrak{b}}
\providecommand{\mx}{\mathfrak{x}}
\providecommand{\my}{\mathfrak{y}}
\providecommand{\mz}{\mathfrak{z}}
\providecommand{\me}{\mathfrak{e}}
\providecommand{\mbarA}{\overline{\mathfrak{A}}}
\providecommand{\mbarB}{\overline{\mathfrak{B}}}
\providecommand{\mbar}[1]{\overline{#1}}
\providecommand{\mC}{\mathfrak{C}}
\providecommand{\mE}{\mathfrak{E}}
\providecommand{\mF}{\mathfrak{F}}
\providecommand{\mT}{\mathfrak{T}}
\providecommand{\mW}{\mathfrak{W}}
\providecommand{\md}{\mathfrak{d}}
\providecommand{\mf}{\mathfrak{f}}
\providecommand{\mm}{\mathfrak{m}}
\providecommand{\mn}{\mathfrak{n}}
\providecommand{\mo}{\mathfrak{o}}
\providecommand{\mpideal}{\mathfrak{p}}
\providecommand{\mq}{\mathfrak{q}}
\providecommand{\mt}{\mathfrak{t}}
\providecommand{\mv}{\mathfrak{v}}
\providecommand{\bara}{\overline{\mathfrak{a}}}



% Batch 25 module/residue macros
\providecommand{\mU}{\mathfrak{U}}
\providecommand{\mP}{\mathfrak{P}}
\providecommand{\mQ}{\mathfrak{Q}}
\providecommand{\mS}{\mathfrak{S}}
\providecommand{\mD}{\mathfrak{D}}
\providecommand{\mR}{\mathfrak{R}}
\providecommand{\mV}{\mathfrak{V}}
\providecommand{\ideal}[1]{\mathfrak{#1}}

% Batch 31 polynomial-ideal/resultant macros
\providecommand{\Amod}{\mathfrak A}
\providecommand{\Bmod}{\mathfrak B}
\providecommand{\Gmod}{\mathfrak G}
\providecommand{\Rmod}{\mathfrak R}
\providecommand{\Smod}{\mathfrak S}
\providecommand{\mideal}{\mathfrak m}
\providecommand{\nideal}{\mathfrak n}
\providecommand{\gideal}{\mathfrak g}
\providecommand{\bb}{\mathfrak b}
\providecommand{\cc}{\mathfrak c}
\providecommand{\zero}[1]{\equiv 0(#1)}
\providecommand{\Norm}{N}


\providecommand{\Amod}{\mathfrak A}
\providecommand{\Bmod}{\mathfrak B}
\providecommand{\Gmod}{\mathfrak G}
\providecommand{\Mmod}{\mathfrak M}
\providecommand{\Nmod}{\mathfrak N}
\providecommand{\Smod}{\mathfrak S}
\providecommand{\mideal}{\mathfrak m}
\providecommand{\nideal}{\mathfrak n}
\providecommand{\gideal}{\mathfrak g}
\providecommand{\hideal}{\mathfrak h}
\providecommand{\jideal}{\mathfrak j}
\providecommand{\tideal}{\mathfrak t}
\providecommand{\aideal}{\mathfrak a}
\providecommand{\bideal}{\mathfrak b}
\providecommand{\zero}[1]{\equiv 0\pmod{#1}}
\providecommand{\Norm}{N}


\providecommand{\frakS}{\mathfrak S}
\providecommand{\frakM}{\mathfrak M}
\providecommand{\frakG}{\mathfrak G}
\providecommand{\frakm}{\mathfrak m}
\providecommand{\frakn}{\mathfrak n}
\providecommand{\frakp}{\mathfrak p}
\providecommand{\frakq}{\mathfrak q}
\providecommand{\frakr}{\mathfrak r}
\providecommand{\fraka}{\mathfrak a}
\providecommand{\frakb}{\mathfrak b}
\providecommand{\frakc}{\mathfrak c}
\providecommand{\frakd}{\mathfrak d}
\providecommand{\frake}{\mathfrak e}
\providecommand{\frakg}{\mathfrak g}
\providecommand{\fraks}{\mathfrak s}
\providecommand{\frako}{\mathfrak o}
\providecommand{\frK}{\mathfrak K}
\providecommand{\barfrakm}{\overline{\mathfrak m}}

\providecommand{\frakh}{\mathfrak h}
\providecommand{\frakt}{\mathfrak t}

% Batch 36 abstract ideal theory / invariant theory macros
\providecommand{\frakR}{\mathfrak R}
\providecommand{\frakT}{\mathfrak T}
\providecommand{\frakH}{\mathfrak H}
\providecommand{\frakP}{\mathfrak P}
\providecommand{\frakQ}{\mathfrak Q}
\providecommand{\frakZ}{\mathfrak Z}
\providecommand{\frakN}{\mathfrak N}
\providecommand{\frakB}{\mathfrak B}
\providecommand{\frakF}{\mathfrak F}
\providecommand{\frakL}{\mathfrak L}
\providecommand{\frakW}{\mathfrak W}
\providecommand{\frakGg}{\mathfrak G}
\providecommand{\Endl}{\operatorname{End}}
\providecommand{\Gal}{\operatorname{Gal}}
\providecommand{\Trdeg}{\operatorname{trdeg}}


% Batch 38 continuation macros
\providecommand{\frR}{\mathfrak R}
\providecommand{\frK}{\mathfrak K}
\providecommand{\frL}{\mathfrak L}
\providecommand{\frM}{\mathfrak M}
\providecommand{\frF}{\mathfrak F}
\providecommand{\frT}{\mathfrak T}
\providecommand{\frN}{\mathfrak N}
\providecommand{\frO}{\mathfrak O}
\providecommand{\frB}{\mathfrak B}
\providecommand{\frC}{\mathfrak C}
\providecommand{\frA}{\mathfrak A}


% Batch 41 Noether macros
\providecommand{\mR}{\mathfrak{R}}
\providecommand{\bR}{\overline{\mathfrak{R}}}
\providecommand{\mK}{\mathfrak{K}}
\providecommand{\mL}{\mathfrak{L}}
\providecommand{\mS}{\mathfrak{S}}
\providecommand{\mT}{\mathfrak{T}}
\providecommand{\mZ}{\mathfrak{Z}}
\providecommand{\mO}{\mathfrak{O}}
\providecommand{\mo}{\mathfrak{o}}
\providecommand{\mD}{\mathfrak{D}}
\providecommand{\mE}{\mathfrak{E}}
\providecommand{\mP}{\mathfrak{P}}
\providecommand{\mf}{\mathfrak{f}}
\providecommand{\mpideal}{\mathfrak{p}}
\providecommand{\mq}{\mathfrak{q}}
\providecommand{\mt}{\mathfrak{t}}
\providecommand{\mOmega}{\mathfrak{\Omega}}
\providecommand{\bP}{\overline{P}}
\providecommand{\bq}{\overline{\mathfrak{q}}}
\providecommand{\bp}{\overline{\mathfrak{p}}}
\providecommand{\ba}{\overline{\mathfrak{a}}}
\providecommand{\bb}{\overline{\mathfrak{b}}}
\providecommand{\tr}{\operatorname{Tr}}
\providecommand{\Norm}{\operatorname{N}}


\providecommand{\rank}{\operatorname{rank}}

% Batch 46 macros
\providecommand{\frakA}{\mathfrak A}
\providecommand{\frakB}{\mathfrak B}
\providecommand{\frakM}{\mathfrak M}
\providecommand{\frakN}{\mathfrak N}
\providecommand{\frako}{\mathfrak o}
\providecommand{\fraka}{\mathfrak a}
\providecommand{\frakc}{\mathfrak c}
\providecommand{\frakp}{\mathfrak p}
\providecommand{\frakO}{\mathfrak O}
\providecommand{\frakP}{\mathfrak P}
\providecommand{\frakD}{\mathfrak D}
\providecommand{\frakG}{\mathfrak G}
\providecommand{\frakK}{\mathfrak K}
\providecommand{\Sp}{\operatorname{Sp}}
\providecommand{\rank}{\operatorname{rank}}
\providecommand{\Trdeg}{\operatorname{Trdeg}}
\providecommand{\charac}{\operatorname{char}}



% Batch 47 macros
\providecommand{\Om}{\Omega}
\providecommand{\Omfrakp}{\Omega_{\mathfrak p}}
\providecommand{\frakA}{\mathfrak A}
\providecommand{\frakB}{\mathfrak B}
\providecommand{\frakC}{\mathfrak C}
\providecommand{\frakG}{\mathfrak G}
\providecommand{\frakH}{\mathfrak H}
\providecommand{\frakP}{\mathfrak P}
\providecommand{\frakp}{\mathfrak p}
\providecommand{\frakq}{\mathfrak q}
\providecommand{\frakr}{\mathfrak r}
\providecommand{\frakS}{\mathfrak S}
\providecommand{\Cl}{\operatorname{Cl}}
\providecommand{\ind}{\operatorname{ind}}
\providecommand{\N}{\operatorname{N}}
\providecommand{\Mat}{\operatorname{M}}
\providecommand{\Gal}{\operatorname{Gal}}
\providecommand{\Norm}{\operatorname{N}}
\providecommand{\nr}{\operatorname{nr}}
\providecommand{\Tr}{\operatorname{Tr}}
\providecommand{\Br}{\operatorname{Br}}


% Batch 48 macros
\providecommand{\frR}{\mathfrak R}
\providecommand{\frS}{\mathfrak S}
\providecommand{\frT}{\mathfrak T}
\providecommand{\frM}{\mathfrak M}
\providecommand{\frN}{\mathfrak N}
\providecommand{\frA}{\mathfrak A}
\providecommand{\frB}{\mathfrak B}
\providecommand{\frX}{\mathfrak X}
\providecommand{\frY}{\mathfrak Y}
\providecommand{\frZ}{\mathfrak Z}
\providecommand{\frG}{\mathfrak G}
\providecommand{\frH}{\mathfrak H}
\providecommand{\frL}{\mathfrak L}
\providecommand{\frU}{\mathfrak U}
\providecommand{\frD}{\mathfrak D}
\providecommand{\frC}{\mathfrak C}
\providecommand{\frP}{\mathfrak P}
\providecommand{\calR}{\mathcal R}
\providecommand{\calS}{\mathcal S}
\providecommand{\calA}{\mathcal A}
\providecommand{\calB}{\mathcal B}
\providecommand{\calT}{\mathcal T}
\providecommand{\End}{\operatorname{End}}
\providecommand{\Aut}{\operatorname{Aut}}
\providecommand{\Mat}{\operatorname{Mat}}
\providecommand{\rad}{\operatorname{rad}}
\providecommand{\rk}{\operatorname{rk}}
\providecommand{\id}{\operatorname{id}}
\providecommand{\op}{\operatorname{op}}


% Batch 49 macros
\providecommand{\Alg}{\mathcal A}
\providecommand{\Abar}{\overline{A}}
\providecommand{\Bbar}{\overline{B}}
\providecommand{\Kbar}{\overline{K}}
\providecommand{\Zbar}{\overline{Z}}
\providecommand{\Gcal}{\mathfrak G}
\providecommand{\Hcal}{\mathfrak H}
\providecommand{\Scal}{\mathfrak S}
\providecommand{\Tcal}{\mathfrak T}
\providecommand{\Zcal}{\mathfrak Z}
\providecommand{\Kcal}{\mathfrak K}
\providecommand{\Pcal}{\mathfrak P}
\providecommand{\rad}{\operatorname{rad}}
\providecommand{\Sp}{\operatorname{Sp}}
\providecommand{\End}{\operatorname{End}}
\providecommand{\Mat}{\operatorname{Mat}}
\providecommand{\Inv}{\operatorname{Inv}}


% Batch 50 macros
\providecommand{\Sp}{\operatorname{Sp}}
\providecommand{\Tr}{\operatorname{Tr}}
\providecommand{\Norm}{\operatorname{N}}
\providecommand{\Mat}{\operatorname{Mat}}
\providecommand{\Gal}{\operatorname{Gal}}
\providecommand{\Cl}{\operatorname{Cl}}
\providecommand{\Gg}{\mathfrak G}
\providecommand{\Hh}{\mathfrak H}
\providecommand{\Ii}{\mathfrak I}
\providecommand{\Jj}{\mathfrak J}
\providecommand{\Aa}{\mathfrak A}
\providecommand{\Bb}{\mathfrak B}
\providecommand{\Cc}{\mathfrak C}
\providecommand{\Oo}{\mathfrak O}
\providecommand{\oo}{\mathfrak o}
\providecommand{\pp}{\mathfrak p}
\providecommand{\PP}{\mathfrak P}
\providecommand{\LL}{\mathfrak L}
\providecommand{\RR}{\mathfrak R}
\providecommand{\ZZ}{\mathfrak Z}
\providecommand{\ee}{\mathfrak e}


% Batch 51 macros
\providecommand{\Sp}{\operatorname{Sp}}
\providecommand{\Trdeg}{\operatorname{Trdeg}}
\providecommand{\frO}{\mathfrak O}
\providecommand{\fro}{\mathfrak o}
\providecommand{\frh}{\mathfrak h}
\providecommand{\frB}{\mathfrak B}
\providecommand{\frA}{\mathfrak A}
\providecommand{\frM}{\mathfrak M}
\providecommand{\frD}{\mathfrak D}
\providecommand{\frd}{\mathfrak d}
\providecommand{\frK}{\mathfrak K}
\providecommand{\frg}{\mathfrak g}
\providecommand{\frp}{\mathfrak p}

% Extra macros collected from the cumulative source for standalone paper builds.
\providecommand{\mat}[1]{\begin{pmatrix}#1\end{pmatrix}}
\providecommand{\Afield}{\mathfrak A}
\providecommand{\Bfield}{\mathfrak B}
\providecommand{\Cfield}{\mathfrak C}
\providecommand{\Sdomain}{\mathfrak S}
\providecommand{\Kfield}{\mathfrak K}
\providecommand{\Rfield}{\mathfrak R}
\providecommand{\Xreal}{\mathfrak X}
\providecommand{\Yreal}{\mathfrak Y}
\providecommand{\Kfield}{\mathfrak{K}}
\providecommand{\Ssys}{\mathfrak{S}}
\providecommand{\Mfield}{\mathfrak{M}}
\providecommand{\tq}{\tau}
\providecommand{\ds}{\displaystyle}
\providecommand{\Lsys}{\mathfrak{L}}
\providecommand{\Jdom}{\mathfrak{J}}
\providecommand{\Gdom}{\mathfrak{G}}
\providecommand{\Omegaint}{[\Omega]}
\providecommand{\Hgrp}{\mathfrak{H}}
\providecommand{\GaloisRes}{\Phi}
\providecommand{\pol}[2]{\left(#1\frac{\partial}{\partial #2}\right)}
\providecommand{\Deltaop}{\Delta\!\left(\frac{\partial}{\partial x},\frac{\partial}{\partial y},\frac{\partial}{\partial z}\right)}
\providecommand{\Omegaop}{\Omega}
\providecommand{\bdelta}{\bar{\delta}}
\providecommand{\tmod}[1]{\;(\operatorname{mod} #1)}
\providecommand{\eqzero}[1]{\equiv 0\,(#1)}
\providecommand{\mJ}{\mathfrak J}
\providecommand{\mS}{\mathfrak S}
\providecommand{\mo}{\mathfrak o}
\providecommand{\ma}{\mathfrak a}
\providecommand{\mb}{\mathfrak b}
\providecommand{\dx}{\dd x}
\providecommand{\dy}{\dd y}
\providecommand{\mc}{\mathfrak{c}}
\providecommand{\OO}{\Omega}
\providecommand{\Th}{\Theta}
\providecommand{\iso}{\simeq}
\providecommand{\normlhd}{\triangleleft}
\providecommand{\mh}{\mathfrak{h}}
\providecommand{\ml}{\mathfrak{l}}
\providecommand{\mr}{\mathfrak{r}}
\providecommand{\ann}{\operatorname{Ann}}
\providecommand{\tuple}[1]{(#1)}
\begin{document}
% Generated from the final cumulative English Noether TeX for reader-facing packaging.
% Source file: Noether_FINAL_Cumulative_English_Translation_MODERN_TEX_numbered_papers_complete_AUDITED.tex
% Paper: 20; source line range: 12376-12611
\section*{20. An Algebraic Criterion for Absolute Irreducibility.}
\begin{center}
By\\[0.25em]
Emmy Noether in Göttingen.\\[0.6em]
Math. Ann. 85 (1922), pp. 26--33
\end{center}

A polynomial -- with coefficients from an arbitrary, abstractly defined field -- is called \emph{absolutely irreducible} if it remains irreducible in the algebraically closed field to which the coefficient domain can be enlarged.\footnote{That every field can be enlarged, essentially uniquely, to an algebraically closed one, i.e. to one in which every polynomial in one variable splits into linear factors, was shown by E. Steinitz in his \emph{Algebraische Theorie der Körper} (J. f. M. 137 (1910), p.~167); he thereby gave the rational equivalent of the fundamental theorem of algebra.} I show in what follows that absolute irreducibility, in contrast with irreducibility relative to a prescribed field, can be grasped \emph{algebraically}; more precisely, the theorem is:

\emph{To every polynomial in $n\ge2$ variables with indeterminate coefficients there can be assigned an integral rational function, with integral coefficients, of these coefficients and of further indeterminates -- the reducibility form -- in such a way that, for every special polynomial of the same degree, the necessary and sufficient condition for absolute irreducibility is the non-vanishing of the reducibility form.} In other words: \emph{for every special system of values of the coefficients for which the reducibility form vanishes identically in the indeterminates and which does not cause a lowering of degree, the special polynomial is reducible, and conversely.}

If, instead of polynomials, one considers homogeneous forms in one more variable, the degree condition is replaced by the simpler condition that the coefficient system not vanish identically.

As a direct consequence of the theorem one also obtains the following, first stated by A. Ostrowski:\footnote{Zur arithmetischen Theorie der algebraischen Größen. Gött. Nachr. 1919, p.~279 (lemma p.~296). -- For the case of fields consisting of ordinary numbers, Ostrowski, as is known to me, also arrived at the main theorem above; he will publish his investigations on this soon. That, in my arrangement of the proof, the main theorem holds for arbitrary, abstractly defined fields rests on the fact that the reducibility form constructed for indeterminates has \emph{integer coefficients independent of the particular field taken as ground field}.}

\emph{Every absolutely irreducible polynomial with algebraic numbers as coefficients remains irreducible modulo every prime ideal of a finite algebraic field $\Omega$ containing the coefficient domain, with the exception of at most finitely many prime ideals of $\Omega$. In particular, $\Omega$ may also be the field of rational numbers.}

I intend to discuss further applications to the theory of relatively integral functions\footnote{D. Hilbert, Mathematische Probleme. Gött. Nachr. 1900, p.~253, Problem 14.} elsewhere.

1. We first show that \emph{every polynomial in $n\ge2$ variables with indeterminate coefficients is absolutely irreducible}. Set
\begin{equation}
\begin{aligned}
F(x)&=\sum A_{i_1\ldots i_n}x_1^{i_1}\cdots x_n^{i_n} &&(i_1+\cdots+i_n\le l),\\
G(x)&=\sum B_{j_1\ldots j_n}x_1^{j_1}\cdots x_n^{j_n} &&(j_1+\cdots+j_n\le m),\\
H(x)&=F(x)G(x)=\sum C_{k_1\ldots k_n}(A,B)x_1^{k_1}\cdots x_n^{k_n} &&(k_1+\cdots+k_n\le l+m),
\end{aligned}
\tag{1}
\end{equation}
where the $A$ and $B$ denote indeterminates, and the $C(A,B)$ become integral bilinear combinations of the $A$ and $B$. The quotients
\[
D=C(A,B)/C_{0\ldots0}(A,B)
\]
are therefore rational functions of the quotients $A/A_{0\ldots0}$ and $B/B_{0\ldots0}$; but the number of these functions is greater than the number of their arguments, since these numbers are respectively
\[
\binom{l+m+n}{n}-1,
\qquad
\binom{l+n}{n}-1,
\qquad
\binom{m+n}{n}-1,
\]
and
\[
\binom{l+m+n}{n}+1>
\binom{l+n}{n}+\binom{m+n}{n}
\quad\text{for } n\ge2,\ l>0,\ m>0.
\]
Thus there is at least \emph{one} algebraic relation among the $D(A,B)$ and consequently also among the $C(A,B)$:
\begin{equation}
 \Phi(Z)\ne0\quad[\text{identically in }Z_{k_1\ldots k_n}],
 \qquad
 \Phi(C(A,B))=0\quad[\text{identically in }A,B],
\tag{2}
\end{equation}
where $\Phi$ denotes an integral polynomial.

A polynomial with \emph{indeterminate $Z$} as coefficients,
\begin{equation}
 E(x)=\sum Z_{k_1\ldots k_n}x_1^{k_1}\cdots x_n^{k_n}
 \qquad(k_1+\cdots+k_n\le l+m),
\tag{3}
\end{equation}
is therefore necessarily \emph{absolutely irreducible} for $n\ge2$.\footnote{For, in the terminology of 2, the indeterminates $Z$ are not zeros of $\mathfrak P$.}

2. The totality of these polynomials $\Phi(Z)$ which vanish identically in $A,B$ under the substitution $Z=C(A,B)$ forms an \emph{ideal of polynomials},\footnote{Such totalities are usually called modules; but in fact their characteristic property -- that with $\Phi_1$ and $\Phi_2$ they contain also $\Phi_1-\Phi_2$, and with $\Phi$ also $\Psi\Phi$, where $\Psi$ is an arbitrary integral polynomial in $Z$ -- is the ideal property.} more precisely a \emph{prime ideal $\mathfrak P$}; for if, under the substitution, a product $\Phi_1\Phi_2$ vanishes, then at least one factor vanishes. Since (2) is satisfied identically in $A,B$, it remains valid for every special system of values $A,B$, where $A,B$ and consequently also $\Gamma=C(A,B)$ denote quantities of some field. \emph{The coefficients $\Gamma$ of every polynomial reducible in an extension field therefore satisfy the conditions $\Phi(\Gamma)=0$; they are zeros of $\mathfrak P$,} or also of the finitely many basis polynomials of $\mathfrak P$. It remains to \emph{prove the converse}.

For this it should be noted that all relations among $A,B,C$ are preserved when, by introducing a new variable $y$, the polynomials (1) are transformed into homogeneous forms $F(x,y)$, $G(x,y)$, $H(x,y)$ of dimensions $l,m,l+m$. Thus coefficients $\Gamma$ also become zeros of $\mathfrak P$ when, for instance, $F(x,y)$\footnote{Polynomials and forms which arise by replacing the indeterminates in $F,G,\ldots,f,g,\ldots$ by special systems of values will throughout be denoted by overlining: $\bar F,\bar G,\ldots,\bar f,\bar g,\ldots$.} becomes a power of $y$; for them the passage to inhomogeneous polynomials causes a lowering of degree of $\bar H(x)$ and not a factorization. It will be shown that, with the exception of this lowering of degree, the converse always holds; in particular, that for \emph{homogeneous forms the converse holds without exception, provided not all $\Gamma$ vanish; in other words, every zero of $\mathfrak P$ is supplied by the parameter representation $\Gamma=C(A,B)$ with quantities $A,B$ from an extension field}.\footnote{That ``in general'' the zeros of $\mathfrak P$ are supplied by the parameter representation follows from the property of the prime ideal $\mathfrak P$ of defining an irreducible algebraic variety. But, as is well known, this does not imply -- something I had originally overlooked and to which A. Ostrowski called my attention some time ago -- that the parameter representation cannot fail for any special system of values. The proof given here rests on more elementary notions.}

I prove this by the direct construction of finitely many functions $\Phi(Z)$: by means of Kronecker's substitution
\[
 x_\lambda=\xi^{d^{\,n-\lambda}}\text{\footnotemark[8]}
\]
\footnotetext[8]{Festschrift, p.~11.}
\setcounter{footnote}{8}
I assign to the polynomials (1) polynomials in \emph{one} variable; I show that gaps occur here in the exponents, and by forming the norm of the corresponding coefficients I arrive at the functions $\Phi(Z)$ and at the desired criterion.

3. Put:
\begin{equation}
\begin{aligned}
 d&=l+m+1; &\qquad i&=i_1d^{n-1}+i_2d^{n-2}+\cdots+i_n;\\
 j&=j_1d^{n-1}+\cdots+j_n; &\qquad k&=k_1d^{n-1}+\cdots+k_n.
\end{aligned}
\tag{4}
\end{equation}
Then, by Kronecker's substitution, the polynomials $F,G,H$ in (1) correspond respectively to the polynomials
\begin{equation}
\begin{aligned}
 f(\xi)&=\sum A_{i_1\ldots i_n}\xi^i &&(i_1+\cdots+i_n\le l),\\
 g(\xi)&=\sum B_{j_1\ldots j_n}\xi^j &&(j_1+\cdots+j_n\le m),\\
 h(\xi)&=\sum C_{k_1\ldots k_n}(A,B)\xi^k &&(k_1+\cdots+k_n\le l+m).
\end{aligned}
\tag{5}
\end{equation}
This correspondence is known to be \emph{one-to-one}, since for all \emph{induced} exponents occurring in (5), $k=0$ implies $k_1=0,\ldots,k_n=0$; and therefore $i+j=i'+j'$ also implies $i_1+j_1=i'_1+j'_1;\ldots;i_n+j_n=i'_n+j'_n$. Hence distinct products $\xi^i\xi^j$ can produce the same exponent $\xi^k$ only when the corresponding products $x_1^{i_1}\cdots x_n^{i_n}x_1^{j_1}\cdots x_n^{j_n}$ also lead to the same exponent. Thus one obtains
\begin{equation}
 h(\xi)=f(\xi)g(\xi),
\tag{6}
\end{equation}
and from the validity of a relation (6), in such a way that in $f(\xi)$ and $g(\xi)$ no exponents other than the induced ones occur, one obtains back again
\begin{equation}
 H(x)=F(x)G(x).
\tag{7}
\end{equation}
Here \emph{gaps actually occur in the induced exponents in $f(\xi)$ and $g(\xi)$}. The highest exponent of $f(\xi)$ is $ld^{n-1}$, while the second highest is $(l-1)d^{n-1}+d^{n-2}$; their difference is $d^{n-2}(d-1)>1$, since $d=l+m+1\ge3$; $n\ge2$; and the same holds for $g(\xi)$.

The relations (6) and (7), valid for the indeterminates $A,B$, remain valid for every special system of values $A,B$. To preserve the degrees, I homogenize (6) and (7) by means of the variables $\eta$ and $y$. \emph{A homogeneous form $H(x,y)$ therefore splits in an algebraic extension field if and only if, in that field, the corresponding binary form $h(\xi,\eta)$ can be split into two factors in such a way that no exponents other than the induced ones occur in the factors.}

4. In order to carry out the formation of the norm mentioned at the end of 2, let
\[
 a_1,b_1;\ldots; a_p,b_p;\quad a_{p+1},b_{p+1};\ldots; a_{p+q},b_{p+q}
\]
be new indeterminates. Set
\begin{equation}
\begin{aligned}
 \varphi(\xi,\eta)&=(a_1\xi+b_1\eta)\cdots(a_p\xi+b_p\eta)
     =r_p\xi^p+r_{p-1}\xi^{p-1}\eta+\cdots+r_0\eta^p,\\
 \psi(\xi,\eta)&=(a_{p+1}\xi+b_{p+1}\eta)\cdots(a_{p+q}\xi+b_{p+q}\eta)
     =s_q\xi^q+s_{q-1}\xi^{q-1}\eta+\cdots+s_0\eta^q,\\
 \chi(\xi,\eta)&=\varphi(\xi,\eta)\psi(\xi,\eta)
     =t_{p+q}\xi^{p+q}+t_{p+q-1}\xi^{p+q-1}\eta+\cdots+t_0\eta^{p+q};
\end{aligned}
\tag{8}
\end{equation}
so that $r,s,t$ denote the homogenized elementary symmetric functions; that is, the symmetric functions of the rows $a,b$ that are linear in each individual row and from which all integral symmetric functions that are homogeneous and of the same dimension in each individual row are built up integrally and with integral coefficients -- and in only one way.

Now form, with further indeterminates $u_{\mu\nu}$, the linear form
\begin{equation}
 \zeta(u,a,b)=\sum u_{\mu\nu}r_\mu s_\nu,
\tag{9}
\end{equation}
where the summation is extended over prescribed indices $\mu$ and $\nu$. The product of $\zeta(u)$ with all its conjugates arising from permutation of the row $a,b$ which are different from $\zeta(u)$ and from one another -- the norm $N(\zeta(u))$ of $\zeta(u)$ with respect to the field of the $t(a,b)$, when $\zeta(u)$ is regarded as a function of the field of the $r_\mu(a,b)$ and $s_\nu(a,b)$ -- is then an integral symmetric function of all rows $a,b$, homogeneous and of the same dimension in each row. Hence, by what was said above, it is a uniquely determined integral function, with integral coefficients, of the $t(a,b)$ and $u_{\mu\nu}$:
\begin{equation}
 N(\zeta(u,a,b))=T(t(a,b),u)\qquad[\text{identically in }a,b,u].
\tag{10}
\end{equation}
The same reasoning shows that also
\[
N(z-\zeta(u))=z^\delta+T_{\delta-1}(t,u)z^{\delta-1}+\cdots+T(t,u)
\qquad[\text{identically in }a,b,u,z]
\]
where all $T$ are integral functions, with integral coefficients, of $t$ and $u$. Both sides of this identity vanish for $z=\zeta(u)$; and indeed, because of the algebraic independence of the elementary symmetric functions $r(a,b)$ and $s(a,b)$, they vanish identically in $r,s$ when $t(a,b)$ is set, according to (8), equal to an integral bilinear combination $w(r,s)$ of the $r$ and $s$, and $\zeta(u)$ is regarded as a function of $r$ and $s$. Thus
\begin{equation}
\zeta(u)^\delta+T_{\delta-1}(w(r,s),u)\zeta(u)^{\delta-1}+\cdots+T(w(r,s),u)=0.
\quad[\text{identically in }r,s,u].\text{\footnotemark[9]}
\tag{11}
\end{equation}
\footnotetext[9]{If the summation in (9) is extended over all indices, (11) represents Kronecker's extension of Gauss's theorem, essentially in Kronecker's arrangement of the proof (Zur Theorie der Formen höherer Stufen, Berliner Ber. 1883, Werke, 2, p.~417).}
\setcounter{footnote}{9}

The identities (10) and (11) remain valid for all special systems of values $\alpha,\beta$ of $a,b$, respectively $\varrho,\sigma$ of $r,s$. If $\varrho,\sigma$ are chosen in particular so that all products $\varrho_\mu\sigma_\nu$ vanish -- where $\mu,\nu$ denote the indices occurring in (9) -- so that $\zeta(u)$ vanishes under the substitution, then, putting $w(\varrho,\sigma)=\tau$, (11) gives
\begin{equation}
 T(\tau,u)=0\qquad[\text{identically in }u].
\tag{12}
\end{equation}

Conversely, let $\tau_0,\ldots,\tau_{p+q}$ be systems of values, not all zero, for which (12) is satisfied. Since in an algebraic extension field there is a decomposition
\begin{equation}
 \bar\chi(\xi,\eta)=\tau_{p+q}\xi^{p+q}+\cdots+\tau_0\eta^{p+q}
=(\alpha_1\xi+\beta_1\eta)\cdots(\alpha_{p+q}\xi+\beta_{p+q}\eta),
\tag{12}
\end{equation}
in which $\alpha$ and $\beta$ do not vanish simultaneously in any factor,\footnote{This rational ``fundamental theorem of algebra'' is the existence theorem on which the exceptionless validity of the converse stated at the end of 2 rests.} although some $\alpha$ or $\beta$ may vanish, we have $\tau=t(\alpha,\beta)$. Substituting these values into (10), using (12), shows that $\zeta(u,\alpha,\beta)$ or one of its conjugate factors vanishes identically in $u$. Thus, after a suitable numbering of $\alpha,\beta$, setting $r(\alpha,\beta)=\varrho$ and $s(\alpha,\beta)=\sigma$ says precisely that $\bar\chi(\xi,\eta)$ admits \emph{at least one decomposition}
\begin{equation}
\bar\chi(\xi,\eta)=\bar\varphi(\xi,\eta)\bar\psi(\xi,\eta)
=\sum \varrho_\kappa\xi^\kappa\eta^{p-\kappa}\cdot\sum \sigma_\lambda\xi^\lambda\eta^{q-\lambda},
\tag{13}
\end{equation}
\emph{in such a way that all products $\varrho_\mu\sigma_\nu$ occurring in $\zeta(u)$ vanish, but not all $\varrho$ or all $\sigma$ vanish.}

5. The degrees $p,q$ in (8) are now taken to be $ld^{n-1}$ and $md^{n-1}$, i.e. the degrees of $f(\xi)$ and $g(\xi)$ in (5). The indices $\mu,\nu$ are split into $\mu^{(1)},\nu^{(1)},\mu^{(2)},\nu^{(2)}$: here $\mu^{(1)}$ runs through all exponents missing from $f(\xi)$, $\nu^{(1)}$ through all exponents of $\xi$ in $\psi(\xi,\eta)$; correspondingly, $\nu^{(2)}$ runs through all exponents missing from $g(\xi)$, and $\mu^{(2)}$ through all exponents of $\xi$ in $\varphi(\xi,\eta)$. Further, let
\[
 e(\xi)=\sum Z_{k_1\ldots k_n}\xi^k\qquad(k_1+\cdots+k_n\le l+m)
\]
be the polynomial in one variable corresponding to the polynomial (3), and let
\[
 S(Z,u)
\]
be the integral polynomial in $Z$ and $u$ obtained from $T(t,u)$ -- the uniquely determined right hand side of (10), where the summation in (9) is extended over the indicated indices $\mu^{(1)},\nu^{(1)},\mu^{(2)},\nu^{(2)}$ -- by replacing the $t$ by the coefficients $Z$ and the corresponding zeros of $e(\xi)$.

If now $r,s$ are replaced by the coefficients $A,B$ and the corresponding zeros of $f(\xi)$ and $g(\xi)$, then by this substitution $\zeta(u)$ \emph{vanishes} for the indicated summation. Moreover $t=w(r,s)$ passes over into the coefficients $C(A,B)$ and the corresponding zeros of $h(\xi)$; hence $T(t,u)$ passes over into $S(C(A,B),u)$. By the identity (11) we therefore get
\begin{equation}
 S(C(A,B),u)=0.\qquad[\text{identically in }A,B,u].
\tag{14}
\end{equation}
Since (14) is preserved for every special system of values, the coefficients $\Gamma=C(A,B)$ of such special $\bar H(x)$, or more generally $\bar H(x,y)$, that split \emph{in an extension field into two factors} satisfy the condition
\begin{equation}
 S(\Gamma,u)=0\qquad[\text{identically in }u].
\tag{15}
\end{equation}

Conversely, suppose that for the coefficients $\Gamma$, not all zero, of a polynomial $\bar H(x)$, respectively form $\bar H(x,y)$, the condition (15) is satisfied. By the above this is equivalent to $T(\tau,u)=0$, where $\tau$ denotes all coefficients of $\bar h(\xi,\eta)$. Hence by (13) there exists in an extension field of the $\Gamma$ a decomposition
\[
 \bar h(\xi,\eta)=\bar f(\xi,\eta)\cdot\bar g(\xi,\eta),
\]
where, because of the prescribed summation, no non-induced exponents occur in $\bar f$ and $\bar g$; consequently there is a relation
\[
 \bar H(x,y)=\bar F(x,y)\cdot\bar G(x,y).
\]
If, for $y=1$, no factor becomes of degree zero -- in particular if the $\Gamma$ do not cause a lowering of the degree of $\bar H(x)$ -- it follows further that
\[
 \bar H(x)=\bar F(x)\cdot\bar G(x).
\]

\emph{Condition (15) is therefore necessary and sufficient for $\bar H(x,y)$, and, when degree lowering is excluded, also for $\bar H(x)$, to split into two factors in an extension field.}

It follows moreover that $S(Z,u)$ cannot vanish identically in $Z$ and $u$, since in 1 it was shown that for $n\ge2$ polynomials with indeterminates as coefficients, hence also the corresponding homogeneous forms in one more variable, are absolutely irreducible. Thus, with $U$ denoting power products of the $u$,
\[
 S(Z,u)=\sum \Phi_\lambda(Z)U_\lambda,
\]
where the $\Phi_\lambda(Z)$ are, by (14), polynomials belonging to $\mathfrak P$. \emph{It is thereby shown that $\mathfrak P$ has no further zeros beyond those given by the parameter representation $\Gamma=C(A,B)$, where $A$ and $B$ belong to an algebraic extension field of the $\Gamma$.} The converse stated at the end of 2 is therefore proved.

6. \emph{The condition of absolute irreducibility} can now be formulated directly. One need only decompose the degree of $E(x)$ in (3), in all possible ways, into two positive summands $l+m$, and form the form $S(Z,u)$ for each such decomposition. The product over these forms,
\begin{equation}
 R(Z,u)=\prod S(Z,u),
\tag{16}
\end{equation}
then forms the \emph{reducibility form of $E(x)$}; for by the above, every \emph{factorization} of $\bar E(x)$, respectively $\bar E(x,y)$, corresponds to the vanishing of one factor of (16), and conversely. Thus the \emph{main theorem stated in the introduction is proved}, and at the same time it is shown how the reducibility form can actually be constructed in finitely many steps.

7. From the \emph{existence of the reducibility form} (16) the \emph{theorem of Ostrowski} mentioned in the introduction follows immediately. Let
\[
 \bar E(x)=\sum \Gamma_{k_1\ldots k_n}x_1^{k_1}\cdots x_n^{k_n}
 \qquad(k_1+\cdots+k_n\le v)
\]
be an absolutely irreducible polynomial with algebraic integers $\Gamma$ as coefficients; let $\mathfrak K$ denote a finite algebraic number field containing all $\Gamma$, and let $\mathfrak p$ be a prime ideal of $\mathfrak K$.

Then the reducibility form $R(Z,u)$ of exact degree for $\bar E(x)$ remains \emph{different from zero} under the substitution $Z=\Gamma$; that is,
\[
 R(\Gamma,u)=\sum P_\lambda U_\lambda\ne0\qquad[\text{identically in }u].
\]
The ideal $\mathfrak r=(P_1,P_2,\ldots)$ is therefore divisible by only a \emph{finite number} of prime ideals of $\mathfrak K$. If $\mathfrak p$ is different from these finitely many, then $R(\Gamma,u)$ does not vanish identically modulo $\mathfrak p$; and since the residue classes modulo $\mathfrak p$ again form a field, it follows that $\bar E(x)$ is \emph{irreducible modulo $\mathfrak p$}, more precisely absolutely irreducible. The same holds for $\bar E(x,y)$, so that no lowering of degree modulo $\mathfrak p$ occurs.

If the coefficients of $\bar E(x)$ are algebraically fractional numbers, then $\mathfrak p$ must also be taken different from the finitely many prime ideals dividing the common denominator $\Delta$; modulo all remaining prime ideals, $\bar E(x)$ may be replaced by $\Delta\cdot\bar E(x)$, so that the preceding hypotheses are satisfied. This proves the theorem.

\begin{flushleft}
Erlangen, August 1921.
\end{flushleft}
\begin{center}
(Received 28. 8. 1921.)
\end{center}

\end{document}
